\chapter{Backward Class Movements}

\section{The Meaning of Backward Class and Common Features}

\begin{KeyConcept}
\begin{itemize}
\item A \textbf{backward class movement} is a movement of the \textbf{non-Brahmins against the Brahmins} (the 'forward' section). The four movements in the syllabus are the \textbf{Satya Shodak Samaj} (Western India/Maharashtra) and the \textbf{Self-Respect Movement, SNDP Movement and Justice Party} (South India) --- largely pre-independence.
\item \textbf{Features common to all backward-class movements}:
\item 1. \textbf{Initiated by non-Brahmins against Brahmins} --- 'non-Brahmin' is a bigger category (Shudras, untouchables, Muslims --- all non-Brahmin), so it was a \textbf{movement of the exploited majority against the exploiter minority}.
\item 2. \textbf{Aim: to restructure society by equitable distribution of the resources usurped by the Brahmins} --- the five resources being \textbf{water, land, education, temple entry and government jobs} (temple entry was a resource because the temple was a site of spirituality, and restricting it maintained Brahmin monopoly).
\item 3. \textbf{Counter-hegemony against cultural hegemony}: the Brahmins owned and controlled the cultural, economic and political system (not just the means of production), so the movement was a counter-hegemony against that cultural hegemony.
\item 4. In Frank Parkin's terms, it was a \textbf{usurpationary social closure} against the Brahmins' \textbf{exclusionary social closure} (their strategies --- scriptures, monopolising education --- to justify their claim over resources).
\item 5. \textbf{Polarisation}: the non-Brahmins consolidated (forgetting their internal caste differences to unite as 'non-Brahmin'), and the Brahmins consolidated in response --- producing a clear-cut polarisation.
\end{itemize}
\end{KeyConcept}

\section{Satya Shodak Samaj (Jyotiba Phule)}

\begin{KeyConcept}
\begin{itemize}
\item \textbf{Jyotiba Phule} (of the \textbf{Mali}/gardener caste, a Shudra) --- humiliated at a Brahmin friend's wedding, he discovered his caste and resolved to \textbf{abolish the caste system} (a radical alteration, like Ambedkar's 'Annihilation of Caste').
\item He criticised the social-religious reformers (Raja Ram Mohan Roy, \textbf{M. G. Ranade}, Dayanand Saraswati, Ishwar Chandra Vidyasagar, Vivekananda) as \textbf{Brahmins operating within a Brahminical framework} --- they cut the shoots (sati, widow remarriage, girl education) but never the \textbf{roots (Brahminism)}.
\item His verdicts: \textbf{'Brahmo Samaj is Brahmin Samaj'; 'Arya Samaj is Aryan Samaj' (Aryan = Brahmin); 'the Indian National Congress is an Upper-Caste Brahmin Congress.'}
\item He organised \textbf{Satya Shodak marriages} (marriage without Brahmins --- no priest, no ritual, a sacred relation between two people only), sought to \textbf{remove Brahmins from birth, death and marriage}, and opened \textbf{schools/education for the lower castes}.
\item \textbf{Reinterpretation of the past}: the Aryan-invasion story is a myth --- the \textbf{Shudras were the original inhabitants and the original kings} (Bali and Sugriva were Shudra kings); the Aryans came from outside and enslaved them. He wrote \textbf{three books}: \textbf{Gulam Giri} (1873, 'slavery' --- equating the Shudra's condition with the Negro slaves of Africa), \textbf{Deen Bandhu} (the thoughts behind the oppressive Aryan Vedic tradition) and \textbf{Shetkariyanch Asud} (how the peasants could improve their lives).
\item \textbf{After Phule's death, Shahu Maharaj} took over the Samaj and \textbf{used it for political gain} --- fighting for \textbf{Kshatriya status} (a place \emph{within} the Varna order, not its abolition) and special political representation --- so the radical reform agenda was sidelined and the movement declined.
\end{itemize}
\end{KeyConcept}

\begin{QuickRevision}
\begin{itemize}
\item Backward class = non-Brahmin (vs Brahmins as forward); movement of exploited majority vs exploiter minority.
\item Common features: non-Brahmin origin; redistribute resources (water, land, education, temple entry, government jobs); counter-hegemony vs cultural hegemony; usurpationary vs exclusionary closure; polarisation.
\item Satya Shodak Samaj (Phule, Mali/Shudra): abolish caste; reformers cut shoots not roots; Brahmo/Arya Samaj = Brahmin/Aryan Samaj; Congress = Brahmin Congress; Satya Shodak marriage; education; reinterpretation (Shudras = original kings; Bali, Sugriva); three books (Gulam Giri 1873, Deen Bandhu, Shetkariyanch Asud); Shahu Maharaj used it politically $\rightarrow$ decline.
\end{itemize}
\end{QuickRevision}

\section{The Tamil Renaissance and the Non-Brahmin Movement}

\begin{KeyConcept}
\begin{itemize}
\item The \textbf{Tamil Renaissance}: classical Tamil texts --- \textbf{Pattu Pattu, Silappadikaram, Manimekalai and the Sangam literature} --- contributed to \textbf{Dravidian consciousness}, just as the Vedas, Puranas, Samhitas and Smritis contributed to Brahmanical consciousness.
\item Tamil scholars believed that the \textbf{Dravidians were the original inhabitants} of the land, and that the Aryans \textbf{imposed Sanskrit culture and the caste system} on them. The rediscovery and reinterpretation of these texts led to the rise of the \textbf{non-Brahmin movement from 1915 onwards}, and to anti-Brahmin resentment against the Brahmin monopoly over resources.
\end{itemize}
\end{KeyConcept}

\section{The Justice Party}

\begin{KeyConcept}
\begin{itemize}
\item The non-Brahmin movement gave rise to the \textbf{Justice Party} (the \textbf{South India Liberation Front}), formed by \textbf{T. M. Nair, Thyagaraja Chetty and Mudaliar}.
\item Its demand (in the \textbf{Non-Brahmin Manifesto}): \textbf{political representation in the councils, and representation in education and government jobs} --- i.e. replacing Brahmins in these positions, rather than fighting Brahminism as a structure.
\item \textbf{Why it collapsed}: its demands (education, jobs, representation) were \textbf{elitist/brahminical demands}, not the demands of the real oppressed (the untouchables and peasants --- whose demands were food, water, freedom from discrimination). Its \textbf{social composition was professional middle class, zamindars and urban business groups} --- so it represented feudal and commercial interests only, and \textbf{confined itself to the quest for government jobs and offices instead of cultural reform or revival of non-Brahmin culture}.
\end{itemize}
\end{KeyConcept}

\begin{QuickRevision}
\begin{itemize}
\item Tamil Renaissance: Pattu Pattu, Silappadikaram, Manimekalai, Sangam literature $\rightarrow$ Dravidian consciousness (vs Brahmanical consciousness from Vedas/Smritis).
\item Dravidians = original inhabitants; Aryans imposed Sanskrit/caste.
\item Justice Party (South India Liberation Front; Nair, Thyagaraja Chetty, Mudaliar): political representation + education + jobs; elitist demands; professional middle class/zamindars; no cultural reform $\rightarrow$ collapsed.
\end{itemize}
\end{QuickRevision}

\section{E. V. Ramaswamy Naicker and the Self-Respect Movement}

\begin{KeyConcept}
\begin{itemize}
\item \textbf{E. V. Ramaswamy Naicker ('Periyar')}: initially a \textbf{Gandhian} (non-violent; supported the \textbf{Vaikom Satyagraha} temple-entry struggle and the Harijan cause). He joined the \textbf{Congress} (which, from 1920, had become inclusive of peasants), but experienced \textbf{caste discrimination at the Erode Congress meeting} (the lunch was arranged by caste) --- so he quit and called the Congress the \textbf{'Brahminical/Upper-Caste Congress'}.
\item Having left both the Justice Party and the Congress, he launched the \textbf{Self-Respect Movement}. He moved from Gandhi's Varnashram belief to an \textbf{Ambedkarite} position --- his objectives were to \textbf{dissolve the Congress and burn the Shastras} (Vedas, Samhitas, Puranas).
\item \textbf{Influenced by Marxism} (adding it to his Dravidian ideology) he attacked the \textbf{political, cultural and economic system}; \textbf{influenced by Phule} he accepted the \textbf{Aryan-invasion theory} and organised \textbf{self-respect (inter-caste) marriages}, doing away with the services of priests in birth, death and marriage.
\item \textbf{Limitations}: (1) its social base was \textbf{confined to the upper non-Brahmin castes}; (2) in \textbf{1944} Periyar \textbf{merged the Justice Party with the Self-Respect League} to form the \textbf{Dravida Kazhagam (Dravidar Kazhagam)} --- turning a social-reform movement into a \textbf{purely political} one (which later became the DMK) --- thereby \textbf{weakening the ideological struggle against Brahmins and Brahminical culture}.
\end{itemize}
\end{KeyConcept}

\begin{QuickRevision}
\begin{itemize}
\item Periyar (E. V. Ramaswamy Naicker): Gandhian $\rightarrow$ Ambedkarite after Erode Congress caste discrimination; called Congress 'Brahminical Congress'.
\item Self-Respect Movement: dissolve Congress + burn Shastras; + Marxism (attack political/cultural/economic system); + Phule (Aryan invasion; self-respect inter-caste marriages).
\item Limitations: upper non-Brahmin base; 1944 merger $\rightarrow$ Dravida Kazhagam $\rightarrow$ political (DMK) $\rightarrow$ weakened anti-Brahmin ideological struggle.
\end{itemize}
\end{QuickRevision}

